% !TeX root = ../机器学习教材（李航）笔记.tex
% 由主文档通过 \input 载入；请编译主文档。
\setcounter{chapter}{2} % 保留原书章号。
\chapter{$k$ 近邻法}
\label{chap:lihang-knn}

给定已经标注类别的样本，一个直接的分类想法是：相近的输入往往具有相近的输出，因此可用新输入附近的已知样本推断它的类别。\emph{$k$ 近邻法}（$k$-nearest neighbor，$k$-NN）把这一想法落实为先找邻居、再作决策的过程。本章依据教材第 3 章，依次讨论近邻分类算法、决定模型的三个基本要素，以及用 kd 树加速近邻搜索的方法；重点区分分类规则本身与实现该规则的搜索算法。近邻法也可用于回归，本章与教材一致，只讨论分类。

\section{$k$ 近邻算法}
\label{sec:lihang-knn-algorithm}

设训练集为 $T=\{(x_i,y_i)\}_{i=1}^{N}$，其中 $N\geq1$，$x_i\in\mathcal X\subseteq\mathbb R^d$，$y_i\in\mathcal Y=\{c_1,\ldots,c_C\}$，$d$ 为特征维数，$C$ 为类别数。给定查询点 $x\in\mathcal X$、距离度量 $\rho$ 和整数 $1\leq k\leq N$，先按 $\rho(x,x_i)$ 从小到大选取 $k$ 个训练样本，再根据它们的标签确定预测类别。为使算法在距离相等时仍有唯一结果，约定距离并列时优先选择索引较小的样本，并把选出的样本索引集合记为 $N_k(x)$；使用索引可以保留坐标相同的不同观测，而不把它们合并成一个点。

\begin{definition}
\label{def:lihang-knn-rule}
在上述约定下，采用\emph{多数表决规则}（majority voting rule）的 $k$ 近邻分类器定义为
\begin{equation}
  \widehat f_T(x)
  =\underset{c_j\in\mathcal Y}{\operatorname{arg\,max}}
    \sum_{i\in N_k(x)}\mathbf 1\{y_i=c_j\},
  \label{eq:lihang-knn-vote}
\end{equation}
其中 $\mathbf 1\{A\}$ 是事件或条件 $A$ 的指示函数，条件成立时取 $1$，否则取 $0$。若多个类别获得相同的最高票数，固定选择类别索引 $j$ 最小者。这里的“多数”指票数最多，不要求某一类别的票数严格超过 $k/2$。
\end{definition}

算法的输入是训练集 $T$、查询点 $x$、近邻数 $k$、距离 $\rho$ 及并列处理规则，输出是 $\widehat f_T(x)$。具体执行时，计算或检索距离最近的 $k$ 个样本，统计这 $k$ 个样本在各类上的票数，再按式~\eqref{eq:lihang-knn-vote} 输出类别。$k=1$ 时得到\emph{最近邻法}（nearest neighbor）：把最近训练样本的类别直接赋给查询点。奇数 $k$ 可以排除二分类中的票数平局，但多分类中即使 $k$ 为奇数也仍可能并列，因此并列规则不能省略。

教材所说的“没有显式的学习过程”，是指基本近邻法不必先拟合一个有限维参数化判别函数；训练样本本身保留了决定预测所需的信息。这并不表示训练阶段完全没有工作：特征变换、$k$ 的选择和搜索索引的建立仍需要数据与计算。分类器随训练集而改变，其预测能力也不能由“记住了训练样本”直接推出。

\section{$k$ 近邻模型的三个基本要素}
\label{sec:lihang-knn-model}

\subsection{由训练样本确定的空间划分}
\label{subsec:lihang-knn-partition}

当训练集、距离、$k$、表决及并列规则都确定后，每个输入都有唯一的预测类别，因此分类器对应于特征空间的一个划分：预测为同一类的点组成一个决策区域。这里的区域一般不是线性子空间；同一类别的区域也可能不连通。近邻法的模型隐含在样本与规则共同形成的划分中，而不必写成一个简单的解析方程。

最近邻法的几何结构最直接。对每个样本点 $x_i$，先定义闭区域
\begin{equation}
  \widetilde V_i
  =\{x\in\mathcal X:\rho(x,x_i)\leq\rho(x,x_j),\ j=1,\ldots,N\}.
  \label{eq:lihang-knn-cell}
\end{equation}
对于互异的样本点和欧氏距离，这些区域是通常的\emph{Voronoi 单元}（Voronoi cell）。它们覆盖特征空间，但可能在等距边界上重叠；把边界按固定的距离并列规则归属后，才得到互不重叠的单元。最近邻法把第 $i$ 个单元中的输入标记为 $y_i$，同类样本的单元可合并为该类别的决策区域。若不同样本具有相同坐标，仍需按索引或另行约定的规则处理，不能仅凭几何距离断言最近邻唯一。

当 $k>1$ 时，先由最近的 $k$ 个样本组合确定局部邻域，再由标签表决决定类别；近邻集合发生变化的边界与最终类别发生变化的边界不一定相同。由此可见，\emph{距离度量、$k$ 值和分类决策规则}共同决定模型，改变其中任何一项都可能改变预测。

\subsection{距离度量}
\label{subsec:lihang-knn-distance}

距离把“相似”转化为可计算的比较，因此也是建模假设的一部分。设 $u=(u^{(1)},\ldots,u^{(d)})^\mathsf T$、$v=(v^{(1)},\ldots,v^{(d)})^\mathsf T\in\mathbb R^d$。对 $1\leq p<\infty$，\emph{$L_p$ 距离}（$L_p$ distance），又称闵可夫斯基距离（Minkowski distance），定义为
\begin{equation}
  L_p(u,v)=\left(\sum_{\ell=1}^{d}|u^{(\ell)}-v^{(\ell)}|^p\right)^{1/p}.
  \label{eq:lihang-knn-lp}
\end{equation}
常用的三个特例是\emph{曼哈顿距离}（Manhattan distance）、\emph{欧氏距离}（Euclidean distance）和 $L_\infty$ 距离：
\begin{align}
  L_1(u,v)&=\sum_{\ell=1}^{d}|u^{(\ell)}-v^{(\ell)}|,
  \label{eq:lihang-knn-l1}\\
  L_2(u,v)&=\left(\sum_{\ell=1}^{d}|u^{(\ell)}-v^{(\ell)}|^2\right)^{1/2},
  \label{eq:lihang-knn-l2}\\
  L_\infty(u,v)&=\max_{1\leq\ell\leq d}|u^{(\ell)}-v^{(\ell)}|.
  \label{eq:lihang-knn-linf}
\end{align}
若 $a=\max_\ell|u^{(\ell)}-v^{(\ell)}|$，则 $a\leq L_p(u,v)\leq d^{1/p}a$，令 $p\to\infty$ 即得到式~\eqref{eq:lihang-knn-linf}。二维空间中，以原点为中心的 $L_1$、$L_2$、$L_\infty$ 单位球分别为菱形、圆形和边平行于坐标轴的正方形，说明不同距离对各方向上的偏离赋予了不同含义。

条件 $p\geq1$ 不宜省略。例如 $0<p<1$ 时，令 $u=(1,0)^\mathsf T$、$v=(0,1)^\mathsf T$、$o=(0,0)^\mathsf T$，有 $L_p(u,v)=2^{1/p}>2=L_p(u,o)+L_p(o,v)$，因而式~\eqref{eq:lihang-knn-lp} 在此范围内不满足距离的三角不等式。另一方面，即使选定欧氏距离，特征的量纲和尺度仍会影响近邻次序；某一坐标的数值尺度很大时，可能主导距离。需要标准化时，应只用训练部分估计变换，再把同一变换用于验证样本与查询点。

\begin{example}
\label{ex:lihang-knn-distance-change}
教材例 3.1 给出 $x_1=(1,1)^\mathsf T$、$x_2=(5,1)^\mathsf T$、$x_3=(4,4)^\mathsf T$。以 $x_1$ 为查询点，在另外两个点中寻找 $L_p$ 距离下的最近邻，并说明结果如何随 $p$ 改变。
\end{example}
\begin{solve}
对于 $1\leq p<\infty$，
\[
  L_p(x_1,x_2)=4,\qquad L_p(x_1,x_3)=3\cdot2^{1/p}.
\]
设 $p_0=\log 2/\log(4/3)\approx2.4094$。比较两式可知：$1\leq p<p_0$ 时最近邻为 $x_2$；$p=p_0$ 时两点等距，按本章约定选择 $x_2$；$p>p_0$ 时最近邻为 $x_3$。特别地，$p=1,2$ 时选 $x_2$，整数 $p\geq3$ 时选 $x_3$；$p=\infty$ 时两距离分别为 $4$ 和 $3$，也选 $x_3$。这说明更换距离会改变局部邻域，因而可能改变同一查询点的预测类别。
\end{solve}

\subsection{$k$ 值的选择}
\label{subsec:lihang-knn-k}

$k$ 控制预测所使用的局部信息范围。较小的 $k$ 使决策更依赖查询点附近的少数样本，能够保留较细的局部结构，却也更容易受噪声或偶然采样波动影响；较大的 $k$ 把更多样本纳入表决，通常减弱局部波动，却可能把距离较远、类别分布不同的样本混在一起，抹平真实的分类边界。教材将其解释为\emph{近似误差}（approximation error）与\emph{估计误差}（estimation error）之间的权衡：局部规则过于粗糙可能造成系统性偏差，局部样本过少则使从有限数据得到的判断不稳定。这里描述的是常见机制，不能理解为任意固定数据集上的两类误差或测试误差都随 $k$ 严格单调变化。

两个极端有助于理解这种权衡。若训练样本坐标互异，$k=1$ 且预测训练样本时允许它自身参与搜索，则每个训练样本都以自身为最近邻，训练误分类率为零；这并不保证对新样本预测准确。若存在坐标相同而标签冲突的样本，上述零训练误差结论也不再自动成立。$k=N$ 时，所有查询点使用同一个训练集表决，预测退化为训练集中的最多数类；即使类别票数并列，固定的表决并列规则也使预测与 $x$ 无关。

实际选择 $k$ 时，可先给定候选集合，再用验证集或交叉验证比较各候选值的误分类率。交叉验证的每一折都只能从该折训练部分中寻找邻居，$k$ 不能超过该部分的样本数；标准化等预处理也应在该训练部分内重新估计。留一验证时应排除被验证样本自身。选定 $k$ 及其他超参数后，再用全部训练数据建立最终分类器，并用未参与选择的测试集评价。所谓“最优 $k$”在这里是给定候选集合与验证方案下的经验选择，不等同于已知总体分布下的理论最优值。

\subsection{多数表决与局部经验风险}
\label{subsec:lihang-knn-decision}

设 $(X,Y)$ 是新样本的随机输入与标签，取分类损失为 $0$--$1$ 损失 $\ell(c,y)=\mathbf 1\{y\ne c\}$。给定训练集 $T$ 后，分类器的总体风险为
\begin{equation}
  R(\widehat f_T)
  =\mathbb E\bigl[\ell(\widehat f_T(X),Y)\bigr]
  =\mathbb P\{Y\ne\widehat f_T(X)\},
  \label{eq:lihang-knn-population-risk}
\end{equation}
期望与概率针对新样本的分布；训练集在此视为固定。这个总体量通常未知。近邻表决直接使用的是查询点附近的样本频数，而不是式~\eqref{eq:lihang-knn-population-risk} 的精确值。

\begin{proposition}
\label{prop:lihang-knn-local-erm}
固定查询点 $x$ 及其近邻索引集 $N_k(x)$。若对这 $k$ 个样本统一预测为某一类别 $c\in\mathcal Y$，则多数表决最小化局部经验风险
\begin{equation}
  \widehat R_x(c)
  =\frac1k\sum_{i\in N_k(x)}\mathbf 1\{y_i\ne c\}.
  \label{eq:lihang-knn-local-risk}
\end{equation}
\end{proposition}
\begin{proof}
对任意 $c\in\mathcal Y$，指示函数恒等式给出
\[
  \widehat R_x(c)
  =1-\frac1k\sum_{i\in N_k(x)}\mathbf 1\{y_i=c\}.
\]
因此最小化 $\widehat R_x(c)$ 等价于最大化类别 $c$ 的票数；票数并列时，各并列类别都达到相同的最小经验风险，再按既定规则选择其一。
\end{proof}

也可以把 $\widehat\eta_j(x)=k^{-1}\sum_{i\in N_k(x)}\mathbf 1\{y_i=c_j\}$ 看作对局部类别比例的估计，表决即选择 $\widehat\eta_j(x)$ 最大的类别。\emph[caution]{局部经验风险最小不等于总体风险最小，也不等于在所有分类函数中最小化整个训练集的经验风险。}要把局部频数解释为对条件类别概率的良好估计，还需要样本与预测任务相匹配，以及所选距离能够反映有意义的局部性；上述代数推导本身没有证明这些统计性质。

\section{$k$ 近邻法的实现：kd 树}
\label{sec:lihang-knn-kdtree}

分类规则确定后，计算上的主要问题是如何找到最近的 $k$ 个样本。\emph{线性扫描}（linear scan）逐一计算查询点到全部样本的距离；对 $L_p$ 距离，计算这些距离需 $O(Nd)$ 次基本运算，再选出最小的 $k$ 个距离。直接全排序需 $O(N\log N)$ 的比较，也可以用大小为 $k$ 的候选堆或选择算法减少筛选开销。线性扫描是明确、可靠的实现基线，其是否足够快取决于数据规模与查询需求，不能笼统断言它不可用。

\emph{kd 树}（kd tree）是用轴对齐切分组织多维样本的二叉树。每个结点对应一个区域，并记录切分轴、切分位置及相关样本；子结点进一步划分该区域。它利用空间距离下界排除不可能包含更优近邻的子树，从而可能减少实际距离计算。名称中的 kd 沿用其多维空间树结构的习惯称呼；本章始终用 $d$ 表示维数、$k$ 表示邻居数，二者不可混淆。\emph{精确搜索时，更换搜索结构不改变分类器，只改变取得同一组近邻的计算过程。}

\subsection{构造平衡 kd 树}
\label{subsec:lihang-knn-kdtree-build}

从包含全部训练样本的根区域开始，每次选择一个坐标轴，用垂直于该轴的超平面切分当前区域。若根结点深度为 $h=0$，循环选择切分轴
\begin{equation}
  \ell(h)=1+(h\bmod d).
  \label{eq:lihang-knn-split-axis}
\end{equation}
于是二维情形依次使用第一、第二、第一坐标轴，三维情形则循环使用三条坐标轴。不同切分轴与切分位置会得到不同的树，但只要搜索与剪枝正确，精确近邻结果应保持一致。

为明确偶数个样本及重复坐标的处理，下面采用与教材例子相容的约定。当前非空索引集为 $M$，$m=|M|$；按第 $\ell$ 个坐标排序，取第 $\lfloor m/2\rfloor+1$ 个样本的该坐标作为切分值 $s$，即偶数个样本时取靠右的中间样本坐标，而不取两个坐标的平均数。把所有满足 $x_i^{(\ell)}=s$ 的样本存放在当前结点，其余样本分别进入左右子树：
\begin{equation}
  B=\{i\in M:x_i^{(\ell)}=s\},\quad
  M_L=\{i\in M:x_i^{(\ell)}<s\},\quad
  M_R=\{i\in M:x_i^{(\ell)}>s\}.
  \label{eq:lihang-knn-kdtree-split}
\end{equation}
$B$ 是当前结点的样本桶；坐标相等的样本都保留各自的索引和标签。

\textbf{构造算法。}输入为带索引的样本坐标、当前索引集 $M$ 和深度 $h$，初始取 $M=\{1,\ldots,N\}$、$h=0$；输出为根结点及递归生成的子树。
\begin{enumerate}
  \item 若 $M$ 为空，返回空树；否则按式~\eqref{eq:lihang-knn-split-axis} 选择切分轴，按上述中间次序统计量确定 $s$。
  \item 建立当前结点，记录切分轴 $\ell$、切分值 $s$ 及样本桶 $B$，并按式~\eqref{eq:lihang-knn-kdtree-split} 分出 $M_L$、$M_R$。
  \item 对 $M_L$、$M_R$ 分别以深度 $h+1$ 递归构造左右子树。若两者均为空，当前结点就是叶结点；完成后返回当前结点。
\end{enumerate}
每个非空结点至少保存一个样本，因此递归必定终止；同一样本恰好保存在一个结点。每个子集的大小不超过 $\lfloor m/2\rfloor$，故树高为 $O(\log N)$。若许多样本在所选坐标上相同，它们可能集中于一个较大的样本桶；树高较小并不意味着每个结点的处理工作都很少。

\begin{example}
\label{ex:lihang-knn-kdtree-build}
教材例 3.2 的二维样本坐标为
\[
  x_1=(2,3)^\mathsf T,\quad x_2=(5,4)^\mathsf T,\quad
  x_3=(9,6)^\mathsf T,\quad x_4=(4,7)^\mathsf T,\quad
  x_5=(8,1)^\mathsf T,\quad x_6=(7,2)^\mathsf T.
\]
按循环坐标轴及上述中间样本坐标约定构造 kd 树。
\end{example}
\begin{solve}
根结点按第一坐标排序，次序为 $2,4,5,7,8,9$，取第 $4$ 个坐标 $7$，故根保存 $x_6=(7,2)^\mathsf T$，切分直线为 $u^{(1)}=7$。左子集为 $\{x_1,x_2,x_4\}$，右子集为 $\{x_5,x_3\}$。左子集按第二坐标排序为 $3,4,7$，保存中间样本 $x_2=(5,4)^\mathsf T$，其左右子结点分别保存 $x_1$、$x_4$。右子集按第二坐标排序为 $1,6$，保存靠右的中间样本 $x_3=(9,6)^\mathsf T$，其左子结点保存 $x_5$，右子树为空。树的父子关系可写为
\[
  \begin{array}{c|c|c}
    \text{当前结点}&\text{左子结点}&\text{右子结点}\\ \hline
    (7,2)^\mathsf T&(5,4)^\mathsf T&(9,6)^\mathsf T\\
    (5,4)^\mathsf T&(2,3)^\mathsf T&(4,7)^\mathsf T\\
    (9,6)^\mathsf T&(8,1)^\mathsf T&\text{空}
  \end{array}
\]
通常统计学定义下，这六个第一坐标的中位数是 $(5+7)/2=6$；教材例子采用的切分值 $7$ 是某个实际样本的坐标。明确使用靠右中间样本，才能同时解释根结点的选择和每个结点保存样本的实现方式。两种中位数约定均可用于设计空间划分，但不能在同一算法描述中混用。
\end{solve}

若每个结点都重新排序当前子集，平衡递归给出的朴素构造时间为 $O(N\log^2 N)$；若用线性时间的次序统计量选择并线性划分，则固定维数下可做到 $O(N\log N)$。这些是具体实现对应的界，而不是仅凭“取中位数”就自动得到的同一个复杂度。保存全部样本及结点记录需要 $O(N)$ 存储；若另存各结点的 $d$ 维包围盒，还需计入相应的维度开销。

\subsection{最近邻搜索、回溯与剪枝}
\label{subsec:lihang-knn-kdtree-search}

沿切分方向下降，只能找到一个局部候选，不能保证真正最近的样本就在第一次进入的叶结点中。正确的搜索必须保留当前已访问样本中的最佳候选，并在回溯时判断其他子树是否仍可能包含更近的点。以下先用欧氏距离说明：若当前候选索引为 $i_\ast$，半径 $r=L_2(x,x_{i_\ast})$，则所有可能改进距离的样本都在以 $x$ 为中心、半径为 $r$ 的开球内；若还要按索引处理等距情形，也应检查球面上的候选。

设某结点用超平面 $u^{(\ell)}=s$ 切分。查询点到该平面的距离为
\begin{equation}
  \Delta=|x^{(\ell)}-s|.
  \label{eq:lihang-knn-plane-bound}
\end{equation}
对于位于查询点另一侧子区域中的任意样本 $u$，都有 $L_2(x,u)\geq|x^{(\ell)}-u^{(\ell)}|\geq\Delta$。因此 $\Delta>r$ 时可跳过另一侧子树；$\Delta\leq r$ 时仍可能存在更优候选，需要继续搜索，或再使用更紧的区域下界判断。条件 $\Delta\leq r$ 只是不能据此剪枝，并不保证其中一定存在更近样本。同一坐标下界对本章的其他 $L_p$ 距离也成立。

\textbf{精确最近邻搜索算法。}输入为非空训练集构造的 kd 树和查询点 $x$，输出为使 $(L_2(x,x_i),i)$ 按字典序最小的索引，即距离最小、等距时索引最小的样本。初始化时无候选，令 $r=+\infty$；从根结点调用下述递归过程。
\begin{enumerate}
  \item 若当前子树为空，立即返回。否则计算当前结点样本桶中各样本到 $x$ 的距离；若距离更小，或距离相同而索引更小，就更新 $i_\ast$ 和 $r$。
  \item 若 $x^{(\ell)}<s$，先递归搜索左子树，否则先递归搜索右子树。把先搜索的一侧称为近侧，另一侧称为远侧；近侧为空时跳过该递归调用，继续后续判断。
  \item 近侧搜索结束后，使用更新后的 $r$ 与式~\eqref{eq:lihang-knn-plane-bound} 比较。若 $\Delta\leq r$，再递归搜索远侧；若 $\Delta>r$，剪去远侧。完成后返回父调用；根调用全部结束时，输出 $i_\ast$ 及其坐标。
\end{enumerate}
这种写法在下降时检查结点样本，在递归返回时完成回溯，与先到达叶结点再向上检查的思路一致，并能直接处理空子树和一个结点保存多个样本的情形。\emph[caution]{某一子树被剪去只表示不再访问这棵子树，仍须返回上层检查其他尚未处理的分支；不能在第一次遇到不相交区域时结束整个搜索。}

\begin{proposition}
\label{prop:lihang-knn-kdtree-correct}
对有限非空训练集，按上述构造与搜索算法执行，最终返回欧氏距离最近、等距时索引最小的训练样本。
\end{proposition}
\begin{proof}
样本桶及左右子树使训练样本被不重不漏地保存。搜索始终把已访问样本中按距离与索引排序的最优者作为当前候选，故 $r$ 在获得第一个候选后只会减小或保持不变。若某子树被剪去，则其中每个样本的距离都大于剪枝当时的 $r$，由于半径不增，也严格大于最终半径，因而不能成为更优者。距离下界等于 $r$ 时没有剪枝，等距且索引更小的候选不会因此遗漏。有限树的递归结束后，每个样本或者已检查，或者处于已证明不能改进结果的子树中，因此当前候选就是全体样本中的最优者。
\end{proof}

若已知某子树样本均包含于轴对齐闭包围盒 $R=\prod_{j=1}^{d}[a_j,b_j]$，还可使用更紧的下界。令
\[
  \delta_j=\max\{a_j-x^{(j)},\ 0,\ x^{(j)}-b_j\},
\]
其中 $\delta_j$ 是查询点第 $j$ 个坐标到区间 $[a_j,b_j]$ 的距离，则
\begin{equation}
  \operatorname{dist}_2(x,R)
  =\inf_{u\in R}L_2(x,u)
  =\left(\sum_{j=1}^{d}\delta_j^2\right)^{1/2}.
  \label{eq:lihang-knn-box-bound}
\end{equation}
等式成立是因为平方距离按坐标相加，每个坐标都可独立投影到对应区间。若这个下界大于 $r$，整棵子树即可剪去；否则继续检查。对 $L_p$ 距离相应使用 $(\sum_j\delta_j^p)^{1/p}$，对 $L_\infty$ 距离使用 $\max_j\delta_j$。剪枝下界必须与实际选用的距离一致。

\begin{example}
\label{ex:lihang-knn-kdtree-query}
在例~\ref{ex:lihang-knn-kdtree-build} 的 kd 树中，按欧氏距离寻找 $x=(3,4.5)^\mathsf T$ 的最近邻。这也是教材习题 3.2 的查询点。
\end{example}
\begin{solve}
先检查根样本 $(7,2)^\mathsf T$，平方距离为 $89/4$。查询点第一坐标小于 $7$，进入左子树；检查 $(5,4)^\mathsf T$ 后，当前最小平方距离变为 $17/4$。查询点第二坐标大于 $4$，先到 $(4,7)^\mathsf T$，其平方距离为 $29/4$，不能改进候选。

回到以第二坐标 $4$ 切分的结点，查询点到切分线的距离为 $|4.5-4|=1/2$，小于当前半径 $\sqrt{17}/2$，故还需检查另一侧的 $(2,3)^\mathsf T$。该点的平方距离为
\[
  (3-2)^2+(4.5-3)^2=\frac{13}{4},
\]
于是更新当前候选为 $(2,3)^\mathsf T$，半径为 $\sqrt{13}/2$。再回到根结点，查询点到切分线 $u^{(1)}=7$ 的距离为 $4>\sqrt{13}/2$，因此右子树可整体剪去。最终最近邻为 $(2,3)^\mathsf T$，距离为 $\sqrt{13}/2$。第一次进入的叶结点保存 $(4,7)^\mathsf T$，而最终结果来自另一分支，说明回溯不可省略。
\end{solve}

从最近邻推广到 $k$ 近邻时，用容量为 $k$ 的\emph{最大堆}（max-heap）保存目前最好的候选，按 $(L_2(x,x_i),i)$ 排序，堆顶是当前候选中最差者。每检查一个样本，若堆未满就加入；若已满且新样本优于堆顶，就替换堆顶。堆未满时令剪枝半径 $r=+\infty$，堆满后令 $r$ 为堆顶距离。近侧优先访问、远侧下界比较和递归返回的过程均沿用上述算法；远侧下界大于 $r$ 时才能剪枝，等于 $r$ 时仍需保留索引并列比较。根调用结束后，堆内即为所需的 $k$ 个近邻，按距离及索引排序即可输出 $N_k(x)$，再交给式~\eqref{eq:lihang-knn-vote} 分类。这给出了教材习题 3.3 所要求的完整推广，$k=1$ 时退化为前述搜索。

最大堆搜索的正确性来自同一不变式：堆中始终保留已访问样本中最优的至多 $k$ 个；堆满后的半径不会增大，剪去的样本也不能进入最终前 $k$ 名。平衡 kd 树的树高虽为 $O(\log N)$，精确查询却可能回溯访问许多分支；在固定低维和有利数据分布下通常能够明显减少访问，但不能仅由“平衡”或“随机分布”断言所有查询都有 $O(\log N)$ 的复杂度。最近邻搜索最坏仍需检查全部样本，距离计算为 $O(Nd)$；维护 $k\geq2$ 个候选时，最坏另有 $O(N\log k)$ 的堆操作。维数较高或空间划分不利时，剪枝可能失效，性能接近线性扫描，因此应结合维数、分布和查询需求选择实现方式。

\section{学习要点与习题研讨}
\label{sec:lihang-knn-review}

本章的逻辑链是：距离确定邻域，$k$ 确定邻域规模，决策规则把局部标签转化为预测；三者与训练集共同形成分类模型。多数表决的局部经验风险解释见命题~\ref{prop:lihang-knn-local-erm}，它不替代对泛化误差的评价。kd 树的空间切分则服务于检索，与最近邻模型的类别区域是两种不同的划分：前者按坐标组织计算，后者由距离和标签决定预测。核查一次 kd 树搜索时，应逐一说明当前候选、半径、待查子树及剪枝下界，而不能只报告下降路径。

\begin{example}
\label{ex:lihang-knn-partition-k}
按教材习题 3.1 的要求自选二维样本，比较 $k=1$ 与 $k=2$ 的决策区域。取 $x_1=(0,0)^\mathsf T$、$y_1=c_1$，$x_2=(2,0)^\mathsf T$、$y_2=c_2$，$x_3=(0,2)^\mathsf T$、$y_3=c_2$，使用欧氏距离，并采用定义~\ref{def:lihang-knn-rule} 的距离和类别并列规则。
\end{example}
\begin{solve}
写查询点为 $x=(u,v)^\mathsf T$，三条平方距离为
\[
  D_1=u^2+v^2,\qquad
  D_2=(u-2)^2+v^2,\qquad
  D_3=u^2+(v-2)^2.
\]
$k=1$ 时，仅 $x_1$ 的类别是 $c_1$，所以预测为 $c_1$ 当且仅当 $D_1\leq D_2$ 且 $D_1\leq D_3$，即 $u\leq1$ 且 $v\leq1$；其余点预测为 $c_2$。等距时索引 $1$ 优先，因此这两条边界在满足另一条件时也归入 $c_1$ 区域。

$k=2$ 时，只要 $x_1$ 位于两个最近邻中，两类便各得一票，按类别索引优先预测为 $c_1$。只有两个近邻恰为 $x_2,x_3$ 时才预测为 $c_2$，这要求 $D_1>D_2$ 且 $D_1>D_3$，即 $u>1$ 且 $v>1$。因此 $k=2$ 时 $c_2$ 区域是这个开象限，其余区域均预测为 $c_1$。若改成票数并列时优先 $c_2$，则 $k=2$ 的预测将处处为 $c_2$，可见并列规则也是模型定义的一部分。

两种 $k$ 的边界不同，但仅凭这三个训练样本及几何划分，无法断定哪种规则对未知分布的新样本更准确；还需独立验证或额外的分布假设。教材习题 3.2 的数值搜索已在例~\ref{ex:lihang-knn-kdtree-query} 完成，习题 3.3 的多近邻搜索则由第~\ref{subsec:lihang-knn-kdtree-search} 小节的候选最大堆给出。
\end{solve}





